\documentclass[10pt,a4paper]{amsart}
\usepackage{amsmath,amssymb,amsthm,graphicx}
% Retained minimal preamble for the Stacks Project chapter flat.tex (More on Flatness).
% Verbatim excerpt of the chapter preamble: Xy-pic setup (native preamble lines 14-18) and the native
% theorem declarations with the mathematical macros (native preamble lines 154-217). Copied mechanically;
% nothing added, removed or redefined. The wrapper class and the amsmath/amssymb/amsthm/graphicx base
% packages are supplied by the assembler, so the class/externaldocument part of the native preamble is omitted.
% For commutative diagrams we use Xy-pic
\usepackage[all]{xy}

% We use 2cell for 2-commutative diagrams.
\xyoption{2cell}
% Theorem environments.
%
\theoremstyle{plain}
\newtheorem{theorem}[subsection]{Theorem}
\newtheorem{proposition}[subsection]{Proposition}
\newtheorem{lemma}[subsection]{Lemma}

\theoremstyle{definition}
\newtheorem{definition}[subsection]{Definition}
\newtheorem{example}[subsection]{Example}
\newtheorem{exercise}[subsection]{Exercise}
\newtheorem{situation}[subsection]{Situation}

\theoremstyle{remark}
\newtheorem{remark}[subsection]{Remark}
\newtheorem{remarks}[subsection]{Remarks}

\numberwithin{equation}{subsection}

% Macros
%
\def\lim{\mathop{\mathrm{lim}}\nolimits}
\def\colim{\mathop{\mathrm{colim}}\nolimits}
\def\Spec{\mathop{\mathrm{Spec}}}
\def\Hom{\mathop{\mathrm{Hom}}\nolimits}
\def\Ext{\mathop{\mathrm{Ext}}\nolimits}
\def\SheafHom{\mathop{\mathcal{H}\!\mathit{om}}\nolimits}
\def\SheafExt{\mathop{\mathcal{E}\!\mathit{xt}}\nolimits}
\def\Sch{\mathit{Sch}}
\def\Mor{\mathop{\mathrm{Mor}}\nolimits}
\def\Ob{\mathop{\mathrm{Ob}}\nolimits}
\def\Sh{\mathop{\mathit{Sh}}\nolimits}
\def\NL{\mathop{N\!L}\nolimits}
\def\CH{\mathop{\mathrm{CH}}\nolimits}
\def\proetale{{pro\text{-}\acute{e}tale}}
\def\etale{{\acute{e}tale}}
\def\QCoh{\mathit{QCoh}}
\def\Ker{\mathop{\mathrm{Ker}}}
\def\Im{\mathop{\mathrm{Im}}}
\def\Coker{\mathop{\mathrm{Coker}}}
\def\Coim{\mathop{\mathrm{Coim}}}

% Boxtimes
%
\DeclareMathSymbol{\boxtimes}{\mathbin}{AMSa}{"02}

%
% Macros for moduli stacks/spaces
%
\def\QCohstack{\mathcal{QC}\!\mathit{oh}}
\def\Cohstack{\mathcal{C}\!\mathit{oh}}
\def\Spacesstack{\mathcal{S}\!\mathit{paces}}
\def\Quotfunctor{\mathrm{Quot}}
\def\Hilbfunctor{\mathrm{Hilb}}
\def\Curvesstack{\mathcal{C}\!\mathit{urves}}
\def\Polarizedstack{\mathcal{P}\!\mathit{olarized}}
\def\Complexesstack{\mathcal{C}\!\mathit{omplexes}}
% \Pic is the operator that assigns to X its picard group, usage \Pic(X)
% \Picardstack_{X/B} denotes the Picard stack of X over B
% \Picardfunctor_{X/B} denotes the Picard functor of X over B
\def\Pic{\mathop{\mathrm{Pic}}\nolimits}
\def\Picardstack{\mathcal{P}\!\mathit{ic}}
\def\Picardfunctor{\mathrm{Pic}}
\def\Deformationcategory{\mathcal{D}\!\mathit{ef}}

% --- appended 2026-10-12 (Chief): the project own definitions of the reference/slogan/history
% display environments, copied mechanically from the shared Stacks preamble lines 8-12
% (\usepackage{verbatim} plus the three \newenvironment lines, which the chapter preamble keeps
% verbatim from the project preamble). Nothing added, removed or redefined: these environments
% carry no mathematical content and the project itself typesets them as comments. Without these
% three lines a counted statement whose own text uses them fails with Environment slogan undefined.
\usepackage{verbatim}
\newenvironment{reference}{\comment}{\endcomment}
\newenvironment{slogan}{\comment}{\endcomment}
\newenvironment{history}{\comment}{\endcomment}
\usepackage{cleveref}
\crefname{originalreference}{original reference}{original references}
\Crefname{originalreference}{Original reference}{Original references}
\makeatletter
\newlabel{situation-free-at-generic-points}{{05MP}{source}{}{sourcepointer-1}{}}
\newlabel{situation-free-at-generic-points@cref}{{[originalreference][0][]05MP}{[1][0][]source}{}{}{}}
\makeatother
\crefname{originalequation}{original equation}{original equations}
\Crefname{originalequation}{Original equation}{Original equations}
\makeatletter
\newlabel{equation-free-at-generic-points}{{05MR}{source}{}{sourcepointer-2}{}}
\newlabel{equation-free-at-generic-points@cref}{{[originalequation][0][]05MR}{[1][0][]source}{}{}{}}
\makeatother
\crefname{originallemma}{original lemma}{original lemmas}
\Crefname{originallemma}{Original lemma}{Original lemmas}
\makeatletter
\newlabel{lemma-free-at-generic-points}{{05MS}{source}{}{sourcepointer-3}{}}
\newlabel{lemma-free-at-generic-points@cref}{{[originallemma][0][]05MS}{[1][0][]source}{}{}{}}
\makeatother
\crefname{originallemma}{original lemma}{original lemmas}
\Crefname{originallemma}{Original lemma}{Original lemmas}
\makeatletter
\newlabel{descent-lemma-descent-data-sheaves}{{02W5}{source}{}{sourcepointer-4}{}}
\newlabel{descent-lemma-descent-data-sheaves@cref}{{[originallemma][0][]02W5}{[1][0][]source}{}{}{}}
\makeatother
\crefname{originallemma}{original lemma}{original lemmas}
\Crefname{originallemma}{Original lemma}{Original lemmas}
\makeatletter
\newlabel{more-morphisms-lemma-separated-locally-quasi-finite-morphisms-fppf-descend}{{02W8}{source}{}{sourcepointer-5}{}}
\newlabel{more-morphisms-lemma-separated-locally-quasi-finite-morphisms-fppf-descend@cref}{{[originallemma][0][]02W8}{[1][0][]source}{}{}{}}
\makeatother
\crefname{originallemma}{original lemma}{original lemmas}
\Crefname{originallemma}{Original lemma}{Original lemmas}
\makeatletter
\newlabel{descent-lemma-descending-fppf-property-immersion}{{02YM}{source}{}{sourcepointer-6}{}}
\newlabel{descent-lemma-descending-fppf-property-immersion@cref}{{[originallemma][0][]02YM}{[1][0][]source}{}{}{}}
\makeatother
\crefname{originallemma}{original lemma}{original lemmas}
\Crefname{originallemma}{Original lemma}{Original lemmas}
\makeatletter
\newlabel{algebra-lemma-generic-flatness-Noetherian}{{051R}{source}{}{sourcepointer-7}{}}
\newlabel{algebra-lemma-generic-flatness-Noetherian@cref}{{[originallemma][0][]051R}{[1][0][]source}{}{}{}}
\makeatother
\crefname{originalequation}{original equation}{original equations}
\Crefname{originalequation}{Original equation}{Original equations}
\makeatletter
\newlabel{equation-iso}{{05MI}{source}{}{sourcepointer-8}{}}
\newlabel{equation-iso@cref}{{[originalequation][0][]05MI}{[1][0][]source}{}{}{}}
\makeatother
\crefname{originallemma}{original lemma}{original lemmas}
\Crefname{originallemma}{Original lemma}{Original lemmas}
\makeatletter
\newlabel{lemma-injectivity-map-source-flat-pure}{{05L8}{source}{}{sourcepointer-9}{}}
\newlabel{lemma-injectivity-map-source-flat-pure@cref}{{[originallemma][0][]05L8}{[1][0][]source}{}{}{}}
\makeatother
\crefname{originallemma}{original lemma}{original lemmas}
\Crefname{originallemma}{Original lemma}{Original lemmas}
\makeatletter
\newlabel{lemma-induction-step-fp}{{05IA}{source}{}{sourcepointer-10}{}}
\newlabel{lemma-induction-step-fp@cref}{{[originallemma][0][]05IA}{[1][0][]source}{}{}{}}
\makeatother
\crefname{originaltheorem}{original theorem}{original theorems}
\Crefname{originaltheorem}{Original theorem}{Original theorems}
\makeatletter
\newlabel{theorem-flattening-map}{{05PF}{source}{}{sourcepointer-11}{}}
\newlabel{theorem-flattening-map@cref}{{[originaltheorem][0][]05PF}{[1][0][]source}{}{}{}}
\makeatother
\crefname{originalreference}{original reference}{original references}
\Crefname{originalreference}{Original reference}{Original references}
\makeatletter
\newlabel{limits-remark-limit-preserving}{{05LX}{source}{}{sourcepointer-12}{}}
\newlabel{limits-remark-limit-preserving@cref}{{[originalreference][0][]05LX}{[1][0][]source}{}{}{}}
\makeatother
\crefname{originalsection}{original section}{original sections}
\Crefname{originalsection}{Original section}{Original sections}
\makeatletter
\newlabel{morphisms-section-constructible}{{054H}{source}{}{sourcepointer-13}{}}
\newlabel{morphisms-section-constructible@cref}{{[originalsection][0][]054H}{[1][0][]source}{}{}{}}
\makeatother
\title{The locus where a quasi-coherent module is free of a given rank at the generic fibre point is representable by a locally closed immersion, of finite presentation when the module is}
\author{The Stacks Project Authors}
\date{stacks-project:flat.tex:local-snapshot-20261011; GNU FDL 1.2 or later}
\begin{document}
\maketitle
\noindent\textit{Editorial note.} All delivered mathematical text is the author's own, copied from the pinned Stacks Project chapter file flat.tex; no statement or proof step was rewritten, reordered or removed.\par\medskip
\noindent\textit{Editorial note.} The preamble is the chapter's own retained minimal preamble (sources/stacks-flat/retained-preamble.tex), a verbatim excerpt of the native Xy-pic setup and of the native theorem declarations and mathematical macros, followed by the project own definitions of the reference, slogan and history display environments (verbatim package plus three newenvironment lines) that the retained file carries at its end and that the project itself typesets as comments. No macro was invented or redefined, and the wrapper class and the amsmath/amssymb/amsthm/graphicx base packages are supplied by the assembler.\par\medskip
\noindent\textit{Editorial note.} The counted statement and proof contain no cross reference, no citation and no included file, so the delivered body is self contained.\par\medskip
\noindent\textit{Editorial note.} Canonical location of the counted proposition: Stacks tag 05UC (sources/stacks-flat/pinned-original-source.tex line 7631).\par\medskip
\noindent\textit{Allowed standard prerequisites (editorial).}\begin{itemize}
\item The functors H\_p and F\_iso and the sheaf property in the fppf topology. Same chapter, Situation labelled situation-free-at-generic-points, equations labelled equation-free-at-generic-points and equation-iso, and Lemma labelled lemma-free-at-generic-points, cited at proof lines 7639-7642 and 7693-7696. Fixes the functors being represented and gives fppf descent for them.
\item Descending representability by a locally closed immersion along an etale or fppf cover. Descent, Lemmas labelled descent-lemma-descent-data-sheaves and descent-lemma-descending-fppf-property-immersion, cited at proof lines 7642-7647; More on Morphisms, Lemma labelled more-morphisms-lemma-separated-locally-quasi-finite-morphisms-fppf-descend, cited at proof lines 7643-7645. Reduces the problem to a Zariski local statement on S.
\item Generic flatness for the fibres and the rank function. Algebra, Lemma labelled algebra-lemma-generic-flatness-Noetherian, cited at proof lines 7657-7659. Supplies that F\_s is locally free near the generic point of the fibre, defining the function p(s).
\item Flattening stratification and the theory of the functor F\_iso. Same chapter, Theorem labelled theorem-flattening-map, cited at proof lines 7698-7700, with Lemmas labelled lemma-injectivity-map-source-flat-pure and lemma-induction-step-fp cited at proof lines 7695-7697. Represents F\_iso by a closed immersion once F is generated by p(s) elements.
\end{itemize}\medskip
\ifdefined\hypertarget\hypertarget{sourcepointer-1}{}\fi
\noindent\textit{Source reference.} Original reference 05MP: sources/stacks-flat/pinned-original-source.tex, source lines 6146--6146. Same-chapter pointer into the author environment situation at source line 6146; invoked at proof line None of the counted proof. It is retained only as a navigation pointer; every load-bearing step is inside the delivered span.\par\medskip
\ifdefined\hypertarget\hypertarget{sourcepointer-2}{}\fi
\noindent\textit{Source reference.} Original equation 05MR: sources/stacks-flat/pinned-original-source.tex, source lines 6166--6166. Same-chapter pointer into the author environment equation at source line 6166; invoked at proof line 7673 of the counted proof. It is retained only as a navigation pointer; every load-bearing step is inside the delivered span.\par\medskip
\ifdefined\hypertarget\hypertarget{sourcepointer-3}{}\fi
\noindent\textit{Source reference.} Original lemma 05MS: sources/stacks-flat/pinned-original-source.tex, source lines 6179--6179. Same-chapter pointer into the author environment lemma at source line 6179; invoked at proof line 7642 of the counted proof. It is retained only as a navigation pointer; every load-bearing step is inside the delivered span.\par\medskip
\ifdefined\hypertarget\hypertarget{sourcepointer-4}{}\fi
\noindent\textit{Source reference.} Original lemma 02W5: sources/stacks-chow/tags.txt, source lines 3792--3792. Cross-chapter standard prerequisite located in the project tag list staged in the local snapshot; invoked at proof line 7644 of the counted proof. It is retained only as a navigation pointer; every load-bearing step is inside the delivered span.\par\medskip
\ifdefined\hypertarget\hypertarget{sourcepointer-5}{}\fi
\noindent\textit{Source reference.} Original lemma 02W8: sources/stacks-chow/tags.txt, source lines 3795--3795. Cross-chapter standard prerequisite located in the project tag list staged in the local snapshot; invoked at proof line 7646 of the counted proof. It is retained only as a navigation pointer; every load-bearing step is inside the delivered span.\par\medskip
\ifdefined\hypertarget\hypertarget{sourcepointer-6}{}\fi
\noindent\textit{Source reference.} Original lemma 02YM: sources/stacks-chow/tags.txt, source lines 3881--3881. Cross-chapter standard prerequisite located in the project tag list staged in the local snapshot; invoked at proof line 7648 of the counted proof. It is retained only as a navigation pointer; every load-bearing step is inside the delivered span.\par\medskip
\ifdefined\hypertarget\hypertarget{sourcepointer-7}{}\fi
\noindent\textit{Source reference.} Original lemma 051R: sources/stacks-chow/tags.txt, source lines 6544--6544. Cross-chapter standard prerequisite located in the project tag list staged in the local snapshot; invoked at proof line 7659 of the counted proof. It is retained only as a navigation pointer; every load-bearing step is inside the delivered span.\par\medskip
\ifdefined\hypertarget\hypertarget{sourcepointer-8}{}\fi
\noindent\textit{Source reference.} Original equation 05MI: sources/stacks-flat/pinned-original-source.tex, source lines 5905--5905. Same-chapter pointer into the author environment equation at source line 5905; invoked at proof line 7691 of the counted proof. It is retained only as a navigation pointer; every load-bearing step is inside the delivered span.\par\medskip
\ifdefined\hypertarget\hypertarget{sourcepointer-9}{}\fi
\noindent\textit{Source reference.} Original lemma 05L8: sources/stacks-flat/pinned-original-source.tex, source lines 5547--5547. Same-chapter pointer into the author environment lemma at source line 5547; invoked at proof line 7692 of the counted proof. It is retained only as a navigation pointer; every load-bearing step is inside the delivered span.\par\medskip
\ifdefined\hypertarget\hypertarget{sourcepointer-10}{}\fi
\noindent\textit{Source reference.} Original lemma 05IA: sources/stacks-flat/pinned-original-source.tex, source lines 3141--3141. Same-chapter pointer into the author environment lemma at source line 3141; invoked at proof line 7694 of the counted proof. It is retained only as a navigation pointer; every load-bearing step is inside the delivered span.\par\medskip
\ifdefined\hypertarget\hypertarget{sourcepointer-11}{}\fi
\noindent\textit{Source reference.} Original theorem 05PF: sources/stacks-flat/pinned-original-source.tex, source lines 6671--6671. Same-chapter pointer into the author environment theorem at source line 6671; invoked at proof line 7697 of the counted proof. It is retained only as a navigation pointer; every load-bearing step is inside the delivered span.\par\medskip
\ifdefined\hypertarget\hypertarget{sourcepointer-12}{}\fi
\noindent\textit{Source reference.} Original reference 05LX: sources/stacks-chow/tags.txt, source lines 7251--7251. Cross-chapter standard prerequisite located in the project tag list staged in the local snapshot; invoked at proof line 7714 of the counted proof. It is retained only as a navigation pointer; every load-bearing step is inside the delivered span.\par\medskip
\ifdefined\hypertarget\hypertarget{sourcepointer-13}{}\fi
\noindent\textit{Source reference.} Original section 054H: sources/stacks-chow/tags.txt, source lines 6640--6640. Cross-chapter standard prerequisite located in the project tag list staged in the local snapshot; invoked at proof line 7721 of the counted proof. It is retained only as a navigation pointer; every load-bearing step is inside the delivered span.\par\medskip
\begin{lemma}
\label{lemma-free-at-generic-points-representable}
In Situation \ref{situation-free-at-generic-points}.
For each $p \geq 0$ the functor $H_p$
(\ref{equation-free-at-generic-points}) is representable
by a locally closed immersion $S_p \to S$. If $\mathcal{F}$
is of finite presentation, then $S_p \to S$ is of finite presentation.
\end{lemma}

\begin{proof}
For each $S$ we will prove the statement for all $p \geq 0$ concurrently.
The functor $H_p$ is a sheaf for the fppf topology by
Lemma \ref{lemma-free-at-generic-points}.
Hence combining
Descent, Lemma \ref{descent-lemma-descent-data-sheaves},
More on Morphisms, Lemma
\ref{more-morphisms-lemma-separated-locally-quasi-finite-morphisms-fppf-descend}
, and
Descent, Lemma \ref{descent-lemma-descending-fppf-property-immersion}
we see that the question is local for the \'etale topology on $S$.
In particular, the question is Zariski local on $S$.

\medskip\noindent
For $s \in S$ denote $\xi_s$ the unique generic point of the fibre $X_s$.
Note that for every $s \in S$ the restriction $\mathcal{F}_s$ of
$\mathcal{F}$ is locally free of some rank $p(s) \geq 0$ in some
neighbourhood of $\xi_s$. (As $X_s$ is irreducible
and smooth this follows from generic flatness for $\mathcal{F}_s$ over
$X_s$, see
Algebra, Lemma \ref{algebra-lemma-generic-flatness-Noetherian}
although this is overkill.) For future reference we note that
$$
p(s) =
\dim_{\kappa(\xi_s)}(
\mathcal{F}_{\xi_s} \otimes_{\mathcal{O}_{X, \xi_s}} \kappa(\xi_s)
).
$$
In particular $H_{p(s)}(s)$ is nonempty and $H_q(s)$ is empty
if $q \not = p(s)$.

\medskip\noindent
Let $U \subset X$ be an open subscheme.
As $f : X \to S$ is smooth, it is open.
It is immediate from (\ref{equation-free-at-generic-points})
that the functor $H_p$ for the pair $(f|_U : U \to f(U), \mathcal{F}|_U)$
and the functor $H_p$ for the pair
$(f|_{f^{-1}(f(U))}, \mathcal{F}|_{f^{-1}(f(U))})$
are the same. Hence to prove the existence of $S_p$ over $f(U)$ we may
always replace $X$ by $U$.

\medskip\noindent
Pick $s \in S$. There exists an affine open neighbourhood $U$
of $\xi_s$ such that $\mathcal{F}|_U$ can be generated by at most
$p(s)$ elements. By the arguments above we see that in order to prove
the statement for $H_{p(s)}$ in an neighbourhood of $s$ we may assume
that $\mathcal{F}$ is generated by $p(s)$ elements, i.e., that there exists
a surjection
$$
u : \mathcal{O}_X^{\oplus p(s)} \longrightarrow \mathcal{F}
$$
In this case it is clear that $H_{p(s)}$ is equal to $F_{iso}$
(\ref{equation-iso}) for the map $u$ (this follows immediately from
Lemma \ref{lemma-injectivity-map-source-flat-pure}
but also from
Lemma \ref{lemma-induction-step-fp}
after shrinking a bit more so that both $S$ and $X$ are affine.)
Thus we may apply
Theorem \ref{theorem-flattening-map}
to see that $H_{p(s)}$ is representable by a closed immersion in a
neighbourhood of $s$.

\medskip\noindent
The result follows formally from the above.
Namely, the arguments above show that locally on $S$ the function
$s \mapsto p(s)$ is bounded. Hence we may use induction
on $p = \max_{s \in S} p(s)$. The functor $H_p$ is representable
by a closed immersion $S_p \to S$ by the above. Replace $S$ by
$S \setminus S_p$ which drops the maximum by at least one and
we win by induction hypothesis.

\medskip\noindent
Assume $\mathcal{F}$ is of finite presentation.
Then $S_p \to S$ is locally of finite presentation by
Lemma \ref{lemma-free-at-generic-points} part (2) combined with
Limits, Remark \ref{limits-remark-limit-preserving}.
Then we redo the induction argument in the paragraph to
see that each $S_p$ is quasi-compact when $S$ is affine:
first if $p = \max_{s \in S} p(s)$, then $S_p \subset S$
is closed (see above) hence quasi-compact. Then $U = S \setminus S_p$
is quasi-compact open in $S$ because $S_p \to S$ is a closed
immersion of finite presentation (see discussion in Morphisms, Section
\ref{morphisms-section-constructible} for example). Then $S_{p - 1} \to U$
is a closed immersion of finite presentation, and so $S_{p - 1}$
is quasi-compact and $U' = S \setminus (S_p \cup S_{p - 1})$
is quasi-compact. And so on.
\end{proof}

\end{document}
